% !TeX root = ../动手学强化学习笔记.tex
% 由主文档通过 \input 载入；请编译主文档。
\setcounter{chapter}{5} % 保留原书章号。
\chapter{动态规划}
\label{ch:rl-dp}

动态规划（dynamic programming，DP）使用已知转移与奖励，把长期决策问题分解为价值备份。它既是一类求解算法，也是检验采样方法的重要基准：在小模型中，学习算法得到的估计可以与精确解直接比较。

\section{迭代策略评估}

固定策略 $\pi$，从任意有限价值表 $V_0$ 开始，反复计算 $V_{k+1}=T_\pi V_k$。由压缩性质，$\|V_k-v_\pi\|_\infty\leq\gamma^k\|V_0-v_\pi\|_\infty$。同步更新先用完整旧表算出新表；原地更新直接覆盖已计算分量，因此同一轮后面的状态可能使用刚得到的新值。

可执行的停止准则是计算当前完整残差 $\delta=\|T_\pi V-V\|_\infty$，当 $\delta\leq(1-\gamma)\epsilon$ 时，当前价值误差至多为 $\epsilon$。同步更新时 $\|V_{k+1}-V_k\|_\infty$ 正好是 $V_k$ 的残差；原地更新的一轮变化则不能不加说明地当作同一个残差。异步更新可以节省计算，但收敛需要所有状态得到充分更新以及对过旧信息的适当控制。

\section{策略改进与策略迭代}

\begin{theorem}[策略改进]
在有限折扣 MDP 中，若 $T_{\pi'}v_\pi\geq v_\pi$ 逐状态成立，则 $v_{\pi'}\geq v_\pi$。特别地，对 $v_\pi$ 贪心的策略满足此条件。
\end{theorem}
\begin{proof}
由单调性，$v_\pi\leq T_{\pi'}v_\pi\leq T_{\pi'}^2v_\pi\leq\cdots$；策略算子迭代收敛到 $v_{\pi'}$，取极限即得结论。贪心时 $T_{\pi'}v_\pi=T_*v_\pi\geq T_\pi v_\pi=v_\pi$。
\end{proof}

策略迭代（policy iteration）交替执行完整策略评估和贪心改进。输入为 $P,r,\gamma$，初始化任意平稳确定性策略；每轮先求 $v_\pi$，再在各状态最大化 $r(s,a)+\gamma\sum_{s'}P(s'\mid s,a)v_\pi(s')$。若原动作仍是最大值动作则保留，否则按固定规则换动作；所有动作均不再改变时输出策略及价值。

在精确计算下，发生严格动作改进就会在至少一个状态严格提高价值；平稳确定性策略只有有限个，且不会回到价值已严格较差的策略，因此算法有限步停止。停止时 $T_*v_\pi=v_\pi$，由唯一性知其最优。近似评估和浮点误差需要容差处理，此有限步论证不直接覆盖任意近似实现。

\section{价值迭代与有限时域逆向递推}

价值迭代（value iteration）直接求解最优算子的不动点：初始化 $V_0$，反复令 $V_{k+1}=T_*V_k$，达到残差阈值后输出当前价值及其贪心策略。它每轮不完成某个策略的精确评估，而是只进行一次最优备份。式\eqref{eq:rl-residual-bound}控制价值误差，上一章的贪心策略界再控制策略损失。

\begin{example}[一状态价值迭代]
某持续任务只有一个状态，两个动作的确定性奖励为 $1$ 和 $2$，执行后均回到原状态。取 $\gamma=0.9,V_0=0$，则 $V_{k+1}=2+0.9V_k$，所以 $V_k=20(1-0.9^k)$，极限为 $20$。价值尚未完全收敛时，贪心策略已选到奖励为 $2$ 的动作，说明价值收敛与策略稳定不必同时发生。
\end{example}

有限时域 $T$ 的问题则从终端值 $V_T(s)=0$ 向前求解：
\[
V_t(s)=\max_a\left[r_t(s,a)+\gamma\sum_{s'}P_t(s'\mid s,a)V_{t+1}(s')\right],\quad t=T-1,\ldots,0.
\]
这是有限次逆向递推，可允许 $\gamma=1$。下标记录剩余决策机会，不能一般地删除。

\section{广义策略迭代与计算代价}

评估使价值靠近当前策略的真实回报，改进使策略偏向当前估计较好的动作。二者相互作用形成广义策略迭代（generalized policy iteration，GPI）。评估可以不完全、状态可以只更新一部分，但获得何种收敛结论仍取决于更新规则，GPI 本身不是对任意交替过程的保证。

若共有 $n$ 个状态、每状态至多 $m$ 个动作，稠密转移下完整最优备份需要 $O(n^2m)$ 运算，存储完整转移表也需相同数量级的空间。稀疏模型按非零后继数计数；策略评估的线性方程可直接求解，但一般稠密求解的计算代价较高。规划计算不消耗新环境交互，两种成本应分别报告。

\section{习题与实验}
\begin{enumerate}
\item 在一状态例子中，求使价值误差小于 $0.1$ 所需的充分迭代次数，并比较策略何时稳定。
\item 给出动作价值并列的模型，说明若每轮任意更换同样好的动作，为什么以“策略完全不变”为停止条件可能产生不必要循环。
\item 构造包含终止状态的小网格，用直接线性求解核验策略评估，再比较策略迭代与价值迭代的输出价值、备份次数及运行时间。
\end{enumerate}





